% =============================================================================
% SYSTEM INSTRUCTION TO LLM:
% This file is a PROMPT TEMPLATE.
% You MUST generate a FULL, FINAL, COMPILABLE LaTeX presentation
% based on this template.
%
% CRITICAL RULES:
% 1. DO NOT change anything in the Preamble.
% 2. You MUST delete all [STYLE REF] frames and the AI INSTRUCTION frame.
% 3. You MUST replace all placeholder texts with real content.
% 4. You MUST preserve the collegeBeamer theme and visual style.
% 5. Keep each slide's body text to roughly 8--10 lines. Page budget: target
%    12--18 slides total, with AT MOST 10 CONTENT slides — the title page, TOC,
%    section-transition pages and the Q&A page are NOT counted in that budget.
% 6. NEVER use \includegraphics. No external image files exist. Draw every chart,
%    diagram and table with tikz / pgfplots / booktabs instead.
% 7. Only these packages are available (already loaded in the preamble): collegebeamer,
%    tikz, pgfplots, amsmath, amssymb, pifont, booktabs, colortbl, multirow.
%    Do NOT add \usepackage lines.
% =============================================================================

\documentclass[aspectratio=169,10pt]{beamer}

% -----------------------------------------------------------------------------
% 1. Preamble & Theme (DO NOT CHANGE)
% -----------------------------------------------------------------------------
\usepackage[zh,bjtu]{collegebeamer}

\usepackage{tikz}
\usetikzlibrary{positioning}
\usepackage{pgfplots}
\pgfplotsset{compat=1.17}
\usepackage{amsmath,amssymb}
\usepackage{pifont}
\usepackage{booktabs}
\usepackage{colortbl}
\usepackage{multirow}

% --- Font Setup ---
% English/Numbers -> Times New Roman

% Disable collegeBeamer auto TOC insertion
\AtBeginSection[]{}

% -----------------------------------------------------------------------------
% 1.1 Title page
% -----------------------------------------------------------------------------
% The theme prints "section number + section name" under every frame title
% automatically (guarded inside collegebeamer.sty, so nothing shows before the
% first \section). No framesubtitle overrides are needed here.

% Title page: empty \subtitle{}/\author{} rows are skipped entirely, so an
% unfilled field never leaves a blank gap between title and date.
\makeatletter
\setbeamertemplate{title page}{
    \hspace{-12mm}
    \begin{beamercolorbox}[wd=0.75\textwidth, sep=10pt, leftskip=8mm]{title}
        {\usebeamerfont{title}\inserttitle\par}
        \vskip4pt\textcolor{maincolor}{\rule{28mm}{1.5pt}}\par
        \ifdefempty{\insertsubtitle}{}{%
            \vskip6pt{\usebeamerfont{subtitle}\insertsubtitle\par}}
        \ifdefempty{\beamer@shortauthor}{}{%
            \vskip6pt{\usebeamerfont{author}\usebeamercolor[fg]{author}\insertauthor\par}}
        \vskip6pt{\usebeamerfont{date}\usebeamercolor[fg]{date}\insertdate\par}
    \end{beamercolorbox}
}
\makeatother

% -----------------------------------------------------------------------------
% 2. Metadata (LLM MUST FILL)
% -----------------------------------------------------------------------------
\title[Short Title]{Full Presentation Title}
% Caller fills these three fields. Leave empty when unknown; never invent placeholders.
\subtitle{}
\author{}
\institute{}
\date{\today}

\begin{document}

% =============================================================================
% FRAME 1: TITLE PAGE
% =============================================================================
\maketitle

% =============================================================================
% FRAME 2: MANUAL TABLE OF CONTENTS
% =============================================================================
\begin{frame}{Outline}
\large\datafont
\renewcommand{\baselinestretch}{1.8}\selectfont
% \makebox keeps a fixed number column so all titles share one left edge

\makebox[2.4em][l]{01} Title of Part One

\makebox[2.4em][l]{02} Title of Part Two

\makebox[2.4em][l]{03} Title of Part Three

\makebox[2.4em][l]{04} Title of Part Four

\makebox[2.4em][l]{05} Summary \& Q\&A
\end{frame}

% =============================================================================
% SECTION 1 (EXAMPLE)
% IMPORTANT: define \section{...} before any content frames of that section
% so the subtitle line is complete (no "0").
% =============================================================================
\section{Title of Part One}

% =============================================================================
% FRAME 3: SECTION TRANSITION (EXAMPLE: SECTION 1)
% LLM MUST generate one such page before EACH section
% =============================================================================
\begin{frame}{Outline}
\large\datafont
\renewcommand{\baselinestretch}{1.8}\selectfont

\textbf{\textcolor{maincolor}{\makebox[2.4em][l]{01} Title of Part One}}

\makebox[2.4em][l]{02} Title of Part Two

\makebox[2.4em][l]{03} Title of Part Three

\makebox[2.4em][l]{04} Title of Part Four

\makebox[2.4em][l]{05} Summary \& Q\&A
\end{frame}

% =============================================================================
% CONTENT FRAME TEMPLATE (LLM MUST USE THIS STRUCTURE)
% =============================================================================
% Content page example 1: Text-left, figure-right layout
\begin{frame}{Content Example: Text Left, Figure Right}
% the theme adds "section number + name" under the title automatically
\small
\begin{columns}[T]
\column{0.55\textwidth}
Write the core analysis content on the left side here.
This area is used for background description, problem definition,
logical reasoning, or methodology explanation.
Text should be concise and conclusion-oriented; avoid vague descriptions.

\vspace{0.6em}
\textbf{Key Takeaway:}
State the central message of this slide in one sentence.

\column{0.40\textwidth}
\centering
% Side panel: FEW categories only (4--6 short labels). Compact width. No xlabel.
% For 12-month / many-label series use the full-width chart on the next example page.
\begin{tikzpicture}
% Boxed axes (no "axis x line/axis y line" keys): pgfplots ignores
% "enlarge x limits" with separate axis lines and clips the outer bars.
\begin{axis}[
    width=5.4cm,
    height=3.6cm,
    ybar,
    bar width=9pt,
    enlarge x limits=0.2,
    ymin=0, ymax=100,
    symbolic x coords={Q1,Q2,Q3,Q4},
    xtick=data,
    ylabel={Sample index},
    tick label style={font=\tiny\datafont},
    label style={font=\scriptsize},
    axis line style={maincolor!50},
    tick style={maincolor!50},
    grid=major,
    grid style={maincolor!12},
]
\addplot[fill=maincolor!40, draw=maincolor] coordinates {
    (Q1,42) (Q2,55) (Q3,61) (Q4,78)
};
\end{axis}
\end{tikzpicture}
\end{columns}
\end{frame}

% Content page example 2: Text-top, figure-bottom layout
\begin{frame}{Content Example: Text Top, Figure Bottom}
\small

% Upper text area
When a series has many categories (for example 12 months), use this
full-width layout below -- not the narrow side-panel chart.
Keep the text concise so the wide chart remains readable.

\vspace{1em}

% Lower figure: MANY categories (12 months) → MUST be wide + rotated x labels
\centering
\begin{tikzpicture}
% Boxed axes here as well: keeps Jan/Dec markers clear of the plot edges.
\begin{axis}[
    width=11.5cm, height=4.2cm,
    enlarge x limits=0.04,
    x tick label style={rotate=45, anchor=east, font=\tiny\datafont},
    ylabel style={font=\scriptsize},
    xlabel style={font=\scriptsize, yshift=-2pt},
    ymin=1900, ymax=2900,
    symbolic x coords={Jan,Feb,Mar,Apr,May,Jun,Jul,Aug,Sep,Oct,Nov,Dec},
    xtick=data,
    ylabel={Sample level},
    xlabel={Month (illustrative)},
    legend style={at={(0.02,0.98)}, anchor=north west, font=\tiny, draw=none, fill=none},
    yticklabel style={font=\tiny\datafont},
    grid=major,
    grid style={maincolor!12},
    axis line style={maincolor!50},
    tick style={maincolor!50},
]
\addplot[
    thick,
    mark=*,
    maincolor,
    mark options={fill=white, draw=maincolor, scale=0.7},
] coordinates {
    (Jan,2040) (Feb,2085) (Mar,2160) (Apr,2330) (May,2345) (Jun,2320)
    (Jul,2420) (Aug,2505) (Sep,2630) (Oct,2740) (Nov,2655) (Dec,2620)
};
\addplot[
    thick,
    dashed,
    maincolor!55,
    mark=square*,
    mark options={fill=white, draw=maincolor!55, scale=0.55},
] coordinates {
    (Jan,2010) (Feb,2050) (Mar,2120) (Apr,2280) (May,2300) (Jun,2290)
    (Jul,2380) (Aug,2460) (Sep,2580) (Oct,2700) (Nov,2620) (Dec,2590)
};
\legend{Series A (sample), Series B (sample)}
\end{axis}
\end{tikzpicture}

\vspace{0.3em}
\centerstate{
    \tiny \emph{Source: data compiled from the source document.}
}
\end{frame}


% Content page example 3: Figure-left, text-right layout
\begin{frame}{Content Example: Figure Left, Text Right}
\small
\begin{columns}[c] % [c] = vertically centered; use [T] for top alignment if text is long

% Left: real TikZ structure diagram (layered pipeline)
\column{0.45\textwidth}
    \centering
    \begin{tikzpicture}[
        font=\scriptsize,
        box/.style={
            draw=maincolor, fill=maincolor!12, rounded corners=2pt,
            minimum height=0.62cm, text width=1.35cm, align=center
        },
        layer title/.style={font=\tiny\bfseries, text=maincolor}
    ]
        % Data layer
        \fill[maincolor!6, rounded corners=3pt] (-0.25,1.05) rectangle (4.85,2.45);
        \node[layer title, anchor=north west] at (-0.15,2.40) {Data layer};
        \node[box] (s1) at (0.7,1.55) {Sensors};
        \node[box] (s2) at (2.5,1.55) {Ingest};
        \node[box] (s3) at (4.2,1.55) {Store};
        % Decision layer
        \fill[maincolor!6, rounded corners=3pt] (-0.25,-0.95) rectangle (4.85,0.70);
        \node[layer title, anchor=north west] at (-0.15,0.65) {Decision layer};
        \node[box] (p1) at (0.7,-0.15) {Feature};
        \node[box] (p2) at (2.5,-0.15) {Model};
        \node[box] (p3) at (4.2,-0.15) {Act};
        % Edges
        \draw[->, thick, maincolor] (s1) -- (s2);
        \draw[->, thick, maincolor] (s2) -- (s3);
        \draw[->, thick, maincolor] (s2) -- (p1);
        \draw[->, thick, maincolor] (p1) -- (p2);
        \draw[->, thick, maincolor] (p2) -- (p3);
        \draw[->, thick, maincolor, dashed]
            (p3.south) -- ++(0,-0.28) -| (p1.south)
            node[pos=0.25, below, font=\tiny, text=maincolor] {feedback};
    \end{tikzpicture}

% Right: text analysis area (50\% width)
\column{0.50\textwidth}
    \textbf{Figure / Object Analysis:}
    
    This section provides a detailed interpretation of the visual element
    on the left. Suitable for presenting:
    \begin{itemize}
        \item Physical structure of experimental apparatus
        \item Specific module details of an algorithm or model
        \item Domain-specific data analysis
    \end{itemize}

    \vspace{1em}
    \textbf{Key Findings:}
    \begin{enumerate}
        \item \ctextbf{Observation 1}: The figure on the left shows\ldots
        \item \ctextbf{Observation 2}: Compared to expectations, the results exhibit\ldots
    \end{enumerate}
\end{columns}
\end{frame}

% =============================================================================
% [STYLE REF] TABLE PAGE (DELETE IN FINAL OUTPUT)
% =============================================================================
\begin{frame}{[STYLE REF] Table Style Example}
\small
\centering
\renewcommand{\arraystretch}{1.15}
\footnotesize
\datafont % lining tabular figures: numerals align and sit on the baseline
% \multirow only in the (uncolored) header rows -- never inside zebra-striped
% body rows, where the row color of the next row would paint over the text.
\rowcolors{3}{maincolor!10}{white}
\begin{tabular}{l c c l}
\toprule
\multirow{2}{*}{Metric} & \multicolumn{2}{c}{Phase} & \multirow{2}{*}{Conclusion} \\
\cmidrule(lr){2-3}
 & A & B & \\
\midrule
Throughput (req/s) & 120 & 158 & Room for improvement \\
Latency (ms) & 340 & 210 & \textcolor{maincolor}{\ding{51}} improved \\
Error rate (\%) & 2.4 & 0.8 & Stable \\
\bottomrule
\end{tabular}

\vspace{1em}
% Block usage (optional): at most one block per slide, for the single takeaway.
\begin{block}{Key message}
State in one sentence the conclusion that the table above supports.
\end{block}
\end{frame}

% =============================================================================
% [STYLE REF] TIKZ FLOW PAGE (DELETE IN FINAL OUTPUT)
% =============================================================================
\begin{frame}{[STYLE REF] Flowchart Style Example}
\centering
\begin{tikzpicture}[
node distance=0.9cm and 0.7cm,
process/.style={rectangle,draw=maincolor,fill=maincolor!10,
minimum height=0.95cm,text width=1.9cm,align=center}
]
\node[process] (a) {Collect};
\node[process,right=of a] (b) {Analyze};
\node[process,right=of b] (c) {Decide};
\node[process,below=of b] (d) {Revise};
\draw[->,thick,maincolor] (a)--(b);
\draw[->,thick,maincolor] (b)--(c);
\draw[->,thick,maincolor] (b)--(d)
    node[midway,right,font=\tiny,text=maincolor] {gap};
\draw[->,thick,maincolor,dashed] (d.west) -| (a.south)
    node[near end,left,font=\tiny,text=maincolor] {retry};
\end{tikzpicture}
\end{frame}

% =============================================================================
% [STYLE REF] TWO-COLUMN ANALYSIS (DELETE IN FINAL OUTPUT)
% =============================================================================
\begin{frame}{[STYLE REF] Two-Column Analysis Example}
\small
% Fill BOTH columns to comparable depth (about 5--6 lines each); never leave
% a column with a single placeholder sentence.
\begin{columns}[T]
\column{0.52\textwidth}
\textbf{Findings (itemize):}
\begin{itemize}
    \item Demand recovered steadily across all four quarters
    \item Unit cost fell as utilisation improved
    \item \textcolor{maincolor}{\ding{51}} Retention stayed above the 90\,\% target
    \item \textcolor{maincolor}{\ding{55}} Overseas segment still lags the plan
    \item Risk exposure is concentrated in two suppliers
\end{itemize}

\column{0.43\textwidth}
\textbf{Action items (enumerate):}
\begin{enumerate}
    \item Lock in supplier pricing before Q3
    \item Re-run the capacity model with new data
    \item Present the revised forecast to the board
    \item Start the overseas pricing pilot
    \item Review retention incentives next quarter
\end{enumerate}
\end{columns}
\end{frame}

% =============================================================================
% AI INSTRUCTION PAGE (DELETE IN FINAL OUTPUT)
% =============================================================================
\begin{frame}{AI GENERATION INSTRUCTIONS}
You MUST:
\begin{itemize}
\item Delete all [STYLE REF] frames and this instruction frame
\item Replace all placeholders with real content
\item Insert a highlighted TOC transition page before each section
\item Define \texttt{\textbackslash section\{...\}} before the first frame of that section
\item Use columns for side-by-side text+figure pages; switch to the full-width
      Example-2 layout whenever a chart has more than about 6 categories
\item Page budget: 12--18 slides total, of which AT MOST 10 are content slides
      (title page, TOC, section transitions and the Q\&A page are not counted)
\item Keep information density high (no empty slides); fill both columns of any
      two-column page to comparable depth
\item Use \texttt{\textbackslash datafont} for numerals in charts and tables
\item Output FULL LaTeX code only
\end{itemize}
\end{frame}

% =============================================================================
% FINAL PAGE
% =============================================================================
\QApage

\end{document}
